% =====================================================================
% Section 2: Related Work
%
% Organized by the dimension that distinguishes EMMET-DP:
%   2.1  Stateless load-aware routing
%   2.2  Backpressure and queue-gradient methods
%   2.3  Learning-based routers
%   2.4  Per-packet state in networking
% =====================================================================

\section{Related Work}
\label{sec:related}

We survey four families of adaptive routing approaches and
position the present work along the dimension that distinguishes
it from each: \emph{whether per-packet trajectory information
participates in the routing decision}.

\subsection{Stateless load-aware routing}
\label{subsec:related-stateless}

The classical line of load-aware routing weights each edge by a
function of its instantaneous utilization and runs a shortest-path
algorithm on the resulting graph. The load-aware shortest-path family (LASP)
multiplies edge latency by $1+\rho(u,v)$, producing a Dijkstra
that prefers under-utilized links. Equal-Cost Multi-Path
(ECMP)~\cite{thaler2000ecmp} hashes flows uniformly across
shortest paths of equal weight, distributing load without explicit
state. Multipath TCP (MPTCP)~\cite{ford2013mptcp} lifts the
multipath decision into the transport layer, allowing a single
connection to use several network paths concurrently with
coupled congestion control.

Within this family, \emph{LASP-aug} (the baseline used throughout
this paper) extends the LASP cost function with a loss-snapshot
term, producing the strongest stateless variant: for every
network state observable at the controller, LASP-aug computes the
cost-minimizing shortest path. EMMET-DP uses identical edge
information, but augments the routing decision with per-packet
mass. The improvement of EMMET-DP over LASP-aug is therefore
attributable specifically to the per-packet state, not to a
richer cost function.

A separate line of work addresses load balancing in the
data-plane fabric of large datacenters and WANs. CONGA and
HULA~\cite{alizadeh2014conga,katta2016hula} make per-flowlet
load-balancing decisions inside the switching fabric using
in-network congestion signals; centralized traffic-engineering
systems such as B4 and SWAN~\cite{jain2013b4,hong2013swan}
plan multipath splits at the controller. These approaches are
complementary rather than competing: they balance load across
flows or flowlets, whereas EMMET-DP modulates the routing
decision for an individual packet based on its trajectory. We
do not include them as experimental baselines because they
operate at a different granularity (flowlet/flow) and on a
different deployment regime (distributed fabric or centralized
TE) than the per-packet admission-time setting we evaluate;
combining EMMET-DP with flowlet-level fabrics is a natural
direction for future work.

None of LASP, ECMP, or MPTCP carries per-packet trajectory
information into subsequent routing decisions. Two packets
arriving at the same node from different histories receive the
same routing recommendation under all three.

\subsection{Backpressure and queue-gradient methods}
\label{subsec:related-backpressure}

Backpressure scheduling~\cite{tassiulas1992stability} and its
delay-aware refinements~\cite{ji2013delay} route packets along
gradients of per-queue backlog, achieving throughput-optimal
behavior under stationary arrivals. A well-known limitation of
pure backpressure is the so-called last-mile inefficiency, in
which packets accumulate near the destination before being
forwarded; subsequent variants combine backpressure with
shortest-path bias to mitigate this effect.

Backpressure is centrally about \emph{network-level state}: the
gradient field is a property of the queues, not of the packets
traversing them. A backpressure router does not distinguish
between two packets at the same node; it forwards both along the
direction of steepest queue descent. This is information-rich at
the network level and information-poor at the packet level. EMMET-DP
operates in the orthogonal direction: it adds packet-level state
on top of a graph-level routing decision.

\subsection{Learning-based routers}
\label{subsec:related-learning}

Q-routing~\cite{boyan1994packet} introduced reinforcement learning
to network routing, with each node maintaining Q-values for
neighbor-conditioned forwarding decisions and updating from
observed delays. More recent work scales this approach with deep
networks: Valadarsky et al.~\cite{valadarsky2017learning} train a
graph-aware policy that achieves competitive performance against
oracle baselines on synthetic topologies, and contemporary surveys
\cite{xiao2020survey} document a growing body of RL routers
targeting datacenter, ISP, and wireless contexts.

Learning-based routers can in principle internalize per-packet
information by including features such as time-in-network,
hop-count-so-far, or origin in the policy input. In practice, most
published policies operate on packet-agnostic state: the input is
the local network situation at the routing node, not the
trajectory of the specific packet under consideration. Where
trajectory features have been used, they appear as additional
inputs to a learned policy rather than as a structural commitment
of the algorithm. EMMET-DP makes the structural commitment: the
mass scalar is part of the algorithm's state, not a feature
optionally fed to a learner.

A second contrast: learning-based routers require training data
and adapt their policy over time. EMMET-DP is a fixed dynamic
program with one tunable parameter; the entire algorithm fits in
roughly $250$ lines of Python. The two approaches are
complementary rather than competing—one could imagine a learned
controller that adjusts $\kappa$ from observed performance, with
EMMET-DP as the inner-loop router.

\subsection{Per-packet state in networking}
\label{subsec:related-perpacket}

The closest line of prior work to EMMET-DP is the family of
mechanisms that introduce per-packet or per-flow state into the
network layer. ECN~\cite{ramakrishnan2001ecn} marks packets that
experienced congestion en route, allowing receivers to signal
congestion back to senders. DCTCP~\cite{alizadeh2010dctcp} uses
this signal to adjust congestion windows proportionally to the
fraction of marked packets, achieving low queue occupancy in
datacenter fabrics. Congestion exposure
(ConEx)~\cite{briscoe2014conex} extends ECN by re-exposing
accumulated congestion at network egress, allowing accountability
across administrative boundaries.

These mechanisms share with EMMET-DP the commitment to carrying
state \emph{about} a packet's experience as it traverses the
network. They differ in how that state is used: ECN, DCTCP, and
ConEx feed back into endpoint congestion-control loops; the
network forwarding decisions remain stateless. EMMET-DP closes
that loop within the network layer: the per-packet state directly
modulates the next routing decision, without round-trip feedback
to an endpoint.

We are aware of no prior work that uses a per-packet scalar to
modulate the cost function of an in-network routing dynamic
program. The combination of (i) trajectory-conditioned per-packet
state, (ii) a structural-rather-than-learned algorithm, and
(iii) measurable improvement on real-network topologies appears
to be novel.
